% =====================================================================
% Project ESCAPE - Final Project Proposal (submission version)
% Team Trimax - CS7.501 Advanced NLP
%
% Self-contained preamble: compiles with a base TeX Live install.
% Body is 4 pages in this layout; references run onto page 5.
%
% TO SUBMIT VIA THE ACL OVERLEAF TEMPLATE: keep everything from
% \begin{document} onward, replace this preamble with
% \documentclass[11pt]{acl}, and carry across the \code macro and the
% two tikz lines below. Drop the \textsection line (local-build only).
% acl.sty's text block is smaller than this one, so expect roughly one
% extra half-page there.
% =====================================================================
\documentclass[10pt,twocolumn,a4paper]{article}
\usepackage[margin=1.9cm,columnsep=0.7cm]{geometry}
\usepackage[utf8]{inputenc}
\usepackage{amsmath,amssymb}
\usepackage{natbib}
\usepackage{url}
\usepackage{booktabs}
\usepackage{array}
\usepackage{enumitem}
\usepackage{caption}
\usepackage{tikz}
\usetikzlibrary{positioning}
\usepackage[colorlinks=true,allcolors=blue]{hyperref}
\bibliographystyle{plainnat}

% Local build only: this TeX install lacks the TS1 fonts OT1 needs for \S.
% The ACL template (Times) has them, so this line can simply be dropped.
\renewcommand{\textsection}{\ensuremath{\mathsection}}

\newcommand{\code}[1]{\texttt{\small #1}}
\captionsetup{font=small,labelfont=bf}
\setlist[itemize]{nosep,leftmargin=3.2mm,label=\textbullet}
\setlist[enumerate]{nosep,leftmargin=5mm}
\setlength{\parskip}{0pt}
\setlength{\abovecaptionskip}{4pt}
\setlength{\belowcaptionskip}{0pt}
\setlength{\textfloatsep}{7pt}
\setlength{\floatsep}{7pt}
\setlength{\intextsep}{7pt}

% Tighter headings than the article defaults.
\makeatletter
\renewcommand\section{\@startsection{section}{1}{\z@}%
  {-1.3ex \@plus -0.4ex \@minus -.2ex}{0.4ex \@plus.2ex}{\normalfont\large\bfseries}}
\renewcommand\subsection{\@startsection{subsection}{2}{\z@}%
  {-1.0ex\@plus -0.3ex \@minus -.2ex}{0.3ex \@plus .2ex}{\normalfont\normalsize\bfseries}}
\renewcommand\paragraph{\@startsection{paragraph}{4}{\z@}%
  {0.6ex \@plus.3ex \@minus.2ex}{-0.6em}{\normalfont\normalsize\bfseries}}
\makeatother

\title{\textbf{Project ESCAPE} \\[1pt]
\large Entropy \& Semantic Compute Allocation in Programming Environments:
Structural Discovery via Entropy in the Byte Latent Transformer
Across Code and Natural Language}

\author{Priyanka Agarwal \small (2024101016) \and Shashwat Mukadam \small (2024102008)
  \and Divyansh Atri \small (2024113001)}
\date{\small Team Trimax \quad$\cdot$\quad CS7.501 Advanced NLP \quad$\cdot$\quad Final Project Proposal}

\begin{document}
\maketitle

\begin{abstract}
\noindent
Language models split text with a vocabulary fixed before training, so boilerplate and
dense logic are chopped alike. The Byte Latent Transformer (BLT)
\citep{pagnoni2024blt} instead reads raw bytes and splits \emph{as it goes}, cutting
wherever its own next-byte prediction becomes uncertain: each cut is the model reporting
where its knowledge ran out. Nobody has checked whether those cuts mean anything. We ask
a narrow, falsifiable version --- \textbf{do BLT's boundaries land at the start of
syntactic constituents in code?}, the Grammatical Branching Hypothesis. We give a
mechanism for why they should, state where it must fail, and test it against a null model
holding boundary density fixed, so a positive result is not an artifact of BLT's threshold
$\tau$. Prose is a built-in negative control: if code and prose do not
separate, we are measuring $\tau$, not grammar. We further test whether an entropy-only
mechanism can be mistaken for semantic risk-awareness in memory-unsafe C++, and show the
two can be told apart. A null result is as informative as a positive one.
\end{abstract}

\section{Introduction}

Tokenization has been the standard front end for language models for a decade, but it
fixes a compression ratio in advance: a tokenizer trained before the model ever sees your
data spends as many tokens on a line of boilerplate as on a line of subtle logic. BLT
removes the fixed vocabulary. A small byte-level language model --- the \emph{patcher} ---
reads raw bytes and measures, at each position, its uncertainty about the next byte, i.e.\
Shannon entropy:
\begin{equation}
H(b_i) = -\sum_v p(v \mid b_{<i}) \log p(v \mid b_{<i})
\label{eq:entropy}
\end{equation}
When $H$ crosses a threshold $\tau$, BLT starts a new \emph{patch}; a large transformer
then processes those variable-length patches, spending more compute where bytes were hard
to predict and less where they were easy.

What matters here is that a boundary is not a preprocessing decision. \textbf{It is a
claim the model makes about its own uncertainty} --- interpretable in a way a BPE merge
table never is, because we can ask whether the claim corresponds to anything independently
verifiable. The BLT authors scale the approach to 8B parameters on 4T training bytes but
report only efficiency and downstream accuracy; they never check whether the induced
patches line up with structure we can check. That is the gap. Code is the right place to
look: it carries machine-checkable ground truth in its abstract syntax tree (AST), and a
grammar with sharper choice points than flexible prose.

\paragraph{Contributions.} (1)~\textbf{A stated mechanism, not a black-box metric}: a
falsifiable account (\S\ref{sec:why}) of \emph{why} boundaries should fall at constituent
starts, including where it must fail. (2)~\textbf{A null model that makes the result mean
something}: a density-matched baseline plus a whitespace baseline (\S\ref{sec:pred}), so
alignment is attributable to grammar, not to boundary density or newlines. (3)~\textbf{A
built-in negative control}: the same mechanism predicts a far weaker effect in prose, so
failure to separate code from prose is itself evidence we are measuring $\tau$.
(4)~\textbf{An honest scope limit}: we test what an entropy-only mechanism \emph{cannot}
see (\S\ref{sec:scope}) rather than overclaiming that surface predictability implies
danger-awareness.

\section{Motivation and Related Work}

Interpretability work on tokenization has focused almost entirely on \emph{static}
vocabularies. Byte-pair encoding \citep{sennrich2016bpe} fixes a merge table before
training, and its failure modes --- poor morphological alignment, unstable segmentation of
rare or multilingual tokens --- are well documented. Tokenization-free models remove the
vocabulary but not the uniformity: CANINE \citep{clark2022canine} and ByT5
\citep{xue2022byt5} spend equal compute on every byte, so they never decide \emph{where} a
boundary belongs. BLT is the first architecture at scale where segmentation is both
byte-level and dynamic --- exactly what makes its boundaries worth probing.

Structural probing on code, such as CodeBERT \citep{feng2020codebert}, and syntax probes on contextual
representations \citep{hewitt2019structural} show that \emph{learned representations} can
encode syntax --- but in every case the model was trained or probed \emph{for} that
structure. None asks whether an unsupervised, compression-only signal, with no
representation-learning objective and no exposure to a parser, recovers syntax as a side
effect. Two gaps follow: whether BLT's patcher implicitly recovers constituent structure,
and whether such recovery is genuine grammar-tracking or merely an artifact of patch
density being pinned by $\tau$ --- a distinction requiring a principled null model, not a
raw overlap score.

\section{Research Objectives}

\begin{enumerate}[label=\textbf{O\arabic*.}]
\item \textbf{Syntactic Alignment.} Do BLT's entropy-triggered boundaries align with AST
node boundaries in C++ and Python, beyond a density-matched baseline?
\item \textbf{Compute Allocation Ratio.} How do mean Bytes-Per-Patch (BPP) and its
variance differ between rigid code syntax and natural language?
\item \textbf{Cross-Language Structural Recognition.} Does patching differ for
indentation-based (Python) versus bracket-based (C++) syntax?
\item \textbf{Robustness to Code Noise.} Do boundaries survive character-level noise
(syntax errors, obfuscation), or does the mechanism fail systemically?
\item \textbf{Semantic vs.\ Surface Compute Allocation.} Is extra compute given to
memory-unsafe C++ operations \emph{fully explained} by the surface unpredictability of
nearby identifiers and casts, or does a residual remain after conditioning on that
confound?
\end{enumerate}

\section{Hypothesis}

\subsection{The Grammatical Branching Hypothesis (H1)}

BLT's patch boundaries are not random compression artifacts: they concentrate at the
\textbf{start} of syntactic constituents rather than inside them, more strongly in code
than in prose, and more strongly at starts than at ends. Precisely, H1 claims an
\emph{average-case asymmetry} between constituent starts and constituent interiors. It
does \textbf{not} claim that starts are points of maximal grammatical choice;
\S\ref{sec:why} explains why that stronger claim is false and why it is not needed. We measure how large the asymmetry is.

\subsection{Why We Expect This to Hold}
\label{sec:why}

The patcher is a small byte-level language model with a short (512-byte) context window
that emits a boundary whenever next-byte entropy crosses $\tau$. A boundary therefore
marks exactly one thing: \textbf{a position where the preceding context stops determining
what comes next.} Everything below follows from that sentence.

\paragraph{Entropy bounds branching; it does not track it.}
It is tempting to read entropy as the local branching factor $B$, the number of
grammatically valid continuations. It is not: entropy depends on the \emph{shape} of the
distribution, not the size of its support. The only guaranteed relation is $H \le \log B$,
and both failure directions are real: ten continuations with one holding $99\%$ of the mass
give $H = 0.08$ nats, far below our $\tau = 1.34$, while four equiprobable ones give
$H = 1.39$ nats and do cross it. Branching is necessary for high entropy but nowhere near
sufficient, and a flat distribution over few options outranks a peaked one over many.

The bound survives in the useful direction. A boundary \emph{certifies} that the local
distribution was genuinely flat: $H > \tau$ implies $B > e^{\tau} \approx 3.8$, so every
entropy-triggered boundary sits where at least four continuations were live. Boundaries
are therefore a \textbf{subset} of high-branching positions --- the subset where the
grammar has not already concentrated its mass on one continuation. Hence a prediction we
state in advance: \textbf{precision against constituent starts should exceed recall.}

\paragraph{The load-bearing claim is about interiors, not starts.}
Our argument nowhere requires constituent starts to be maximal-choice points, only that
entropy be lower \emph{inside} a committed production than at its boundary --- a
structural property of context-free grammars, not a statistical hope. Once a production is
entered its own rule constrains what may follow: mid-keyword and mid-identifier bytes are
cheap, and closing delimiters (\code{;}, \code{)}, \code{\}}) are among the most
predictable bytes in a file because the grammar \emph{requires} them once the production is
committed --- they are not chosen. This is the asymmetry P2 rests on: boundaries are pushed
\emph{away from} interiors and right edges, and constituent starts inherit them by
elimination rather than by being choice-maximal.

\paragraph{Where this mechanism does not hold.}
Two dissociations follow; they bound what a positive result can mean, and they damage
\emph{different} metrics. First, some constituents open deterministically: in both C++ and
Python, \code{if}, \code{for} and \code{while} take a forced next token, and \code{def}
must be followed by a name. These are constituent starts with branching factor near one,
and the mechanism predicts BLT will miss them --- a \textbf{recall} cost, tested directly
by P4. Second, high entropy arises where no constituent opens at all: partway through an
identifier, inside a string literal, inside a comment. These are lexical choice points, not
grammatical ones --- a \textbf{precision} cost, and the confound isolated in
\S\ref{sec:scope}. The patcher measures next-byte predictability, conflating grammatical
branching with lexical arbitrariness, and nothing in the architecture separates them.
Table~\ref{tab:example} walks one line of C++ through all three cases.

\begin{table}[t]
\small\centering
\begin{tabular}{@{}llcl@{}}
\toprule
Position & Entropy & Split? & Case \\
\midrule
\code{i} of \code{if}      & high & yes & correct \\
\code{(} after \code{if}   & low  & no  & \emph{missed} start \\
\code{x} of \code{x > 0}   & high & yes & lexical, not syntactic \\
\code{)} \code{;} \code{\}} & low & no  & correct (supports P2) \\
\bottomrule
\end{tabular}
\caption{The mechanism applied to \code{if (x > 0) \{ total += 1; \}}. Row 1 is the effect
we predict; row 2 costs recall (forced opener); row 3 costs precision (identifier entropy);
row 4 is why starts should beat ends.}
\label{tab:example}
\end{table}

\paragraph{Prose is the contrast case.}
Branching is flatter across a sentence: content words are hard everywhere, function words
easy everywhere, and no formal rule forces a clause to end with a specific character. Prose
should therefore show more uniform patch lengths and weaker structural anchoring.
\textbf{This code/prose gap is what makes the claim one about grammar rather than about
$\tau$.}

\subsection{Predictions and Null Model}
\label{sec:pred}

\paragraph{P1: Alignment.} Boundary positions match AST node start offsets, within a
tolerance $k$ fixed \emph{a priori} per language on a held-out calibration split, well
above baseline, for statement- and expression-level nodes.

\paragraph{P2: Asymmetry.} Alignment at node starts is clearly higher than at node ends:
closing delimiters are low-entropy, so BLT should systematically miss right edges. Span-IoU
therefore \emph{understates} the effect, so boundary precision/recall against start offsets
--- not IoU --- is our primary metric.

\paragraph{P3: Cross-domain contrast.} BPP variance is high in code and low in prose, and
P1-style alignment in prose (against a constituency parse) is near chance. Code and prose
must dissociate.

\paragraph{P4: Stratification by opener.} Alignment is measurably weaker for constituent
types with deterministic openers (\code{if\_statement}, \code{for\_statement}) than for
open-ended ones (\code{expression\_statement}, \code{return\_statement}, assignment
right-hand sides). P4 follows from \S\ref{sec:why} and is the sharpest test that the
mechanism running is ours: a generic ``entropy correlates with structure'' story does not
predict it.

\paragraph{Null model (H0).} Agreement between patch boundaries and AST node starts is no
better than a density-matched random boundary set: the same number of boundaries, patch
lengths drawn from BLT's own empirical patch-length distribution, in shuffled order. This
baseline is required, not optional --- mean patch length is roughly pinned by $\tau$, so
\emph{any} boundary set of the right density scores nonzero raw alignment, and only the
margin over H0 is informative. We assess it by permutation test (10{,}000 resamples),
reporting effect size with confidence intervals rather than a raw overlap number.

\paragraph{Second baseline: whitespace.} H0 is structure-blind by construction, so it cannot
rule out the most obvious confound in code: most AST node starts sit immediately after a
newline and indentation, so a trivial ``split at every newline'' rule would score well on
P1. It is therefore an explicit second baseline: H1 is only interesting if BLT beats \emph{that},
not merely chance. This matters most in Python, where indentation \emph{is} the
grammar.

\paragraph{Depth control.} Alignment is additionally computed \emph{per AST depth}: every
byte sits inside several nested nodes, so ``alignment'' is undefined until granularity is
fixed.

\paragraph{What counts as a result.} H1 is supported if P1 clears both baselines with
non-overlapping confidence intervals, P2 holds in the predicted direction, P3 shows
code/prose dissociation, and P4's stratification appears. Any subset failing is reported as
such, and each failure mode is itself interpretable.

\subsection{Why the Result Is Not Obvious in Advance}
\label{sec:notobvious}

If H1 holds, an unsupervised, entropy-driven compressor recovers syntactic constituent
boundaries with no parser in the loop --- an interpretability result about what dynamic
patching has implicitly learned, and evidence that BLT transfers to code without a
code-specific tokenizer.

If H1 fails, \emph{how} it fails is the signal. Chance alignment with uniformly
\textbf{high} entropy means the patcher is undertrained on code and allocates compute blind
to grammar --- a reportable limitation with a clear fix: retrain the entropy model on code.
Chance alignment with uniformly \textbf{low} entropy means code is so compressible that
$\tau$ rarely fires, making boundaries density-limited rather than structure-driven. H1
could also lose outright if identifier entropy dominates: names are arbitrary strings, not
grammar-constrained, and could swamp the syntactic signal --- the precision failure above,
addressed in \S\ref{sec:scope}. Specifying distinct, diagnosable failure modes rather than
a single pass/fail outcome is what makes the experiment informative either way.

\subsection{Scope Limit: What Entropy Cannot See}
\label{sec:scope}

The patcher is a short-context, surface-statistics model with no representation of pointer
aliasing, ownership, or memory safety. Tokens like \code{new}, \code{delete}, \code{*} and
\code{\&} are short, frequent and lexically predictable, so \S\ref{sec:why} predicts
\textbf{low} entropy at them, not high. Consequently \textbf{we do not expect BLT to
allocate compute by semantic risk} --- that claim would contradict the mechanism our own
hypothesis relies on. Any extra compute in memory-unsafe regions should instead come from
the arbitrary identifiers and casts that surround pointer operations. Objective~5 tests this
directly: we condition memory-specific compute profiling on identifier-token entropy via
regression, and predict any apparent memory-safety effect shrinks toward zero once that
confound is controlled. This is a confound-isolation test, not a search for semantic
understanding; a null result here is informative, not a failed experiment.

\begin{figure*}[t]
\centering
\begin{tikzpicture}[scale=0.72,transform shape,
  box/.style={draw,rounded corners=2pt,align=center,inner sep=3pt,font=\footnotesize,
              minimum height=7mm},
  >=stealth]
\node[box] (src) {Source file\\(raw UTF-8 bytes)};
\node[box,right=9mm of src] (patch) {BLT patcher\\\scriptsize 14L, 768d, 512-byte ctx};
\node[box,right=8mm of patch] (ent) {Per-byte\\entropy $H(b_i)$};
\node[box,right=8mm of ent] (bnd) {Boundaries\\\scriptsize $H > \tau$, entropy-triggered};
\node[box,below=7mm of patch] (ts) {\code{tree-sitter}\\parse};
\node[box,below=7mm of ent] (ast) {AST node\\start/end offsets};
\node[box,below=7mm of bnd] (base) {Baselines\\\scriptsize H0 density-matched\\\scriptsize + whitespace};
\node[box,right=9mm of bnd,yshift=-7mm,fill=black!5] (score)
  {\textbf{Boundary P/R}\\\scriptsize at constituent starts\\\scriptsize by depth, by node type};
\draw[->] (src) -- (patch);
\draw[->] (patch) -- (ent);
\draw[->] (ent) -- (bnd);
\draw[->] (src.south) |- (ts.west);
\draw[->] (ts) -- (ast);
\draw[->] (ast) -- (base);
\draw[->] (bnd.east) -| (score.north);
\draw[->] (base.east) -| (score.south);
\end{tikzpicture}
\caption{Extraction and scoring pipeline. The upper path is unsupervised --- the patcher never
sees a parse tree; the lower path supplies independently verifiable ground truth. Boundaries are
scored against AST node starts and against both baselines at once, stratified by AST depth
(\S\ref{sec:pred}) and node type (P4).}
\label{fig:pipeline}
\end{figure*}

\section{Methodology}

Figure~\ref{fig:pipeline} shows the pipeline end to end.

\paragraph{Model and data.} We use \code{BltForCausalLM} from Hugging Face
\code{transformers} with the \code{itazap/blt-1b-hf} checkpoint (a 1B port of BLT's Local
Encoder $\rightarrow$ Latent Transformer $\rightarrow$ Local Decoder architecture, the largest
openly available BLT checkpoint). Natural language comes from
WikiText-103 and code from CodeSearchNet, filtered for Python and C++; for the prose
negative control (P3) we additionally parse WikiText-103 sentences with a standard
constituency parser to obtain phrase-boundary ground truth.

\paragraph{Boundary extraction.} BLT resolves patches dynamically from the entropy threshold
($\tau \approx 1.34$) and an optional maximum patch length. The patcher exposes per-byte
entropies and patch lengths directly, so boundaries are read from its output, not by instrumenting
internals. Each is tagged \emph{entropy-triggered} or \emph{length-capped}: a
patch hitting the length cap before entropy crosses $\tau$ has nothing to do with grammar,
so length-capped boundaries are excluded from the primary P1/P2 scores. Entropy is reported
in nats, so $\tau = 1.34$ nats $= 1.93$ bits.

\paragraph{Checkpoint fidelity.} The only public checkpoint is a third-party conversion, so
we sanity-check it before producing any numbers: coherent generation, sensible boundaries on
simple inputs, and --- critically --- that the patcher's effective context window and
patch-length distribution match the original work. The short
context window is load-bearing for \S\ref{sec:why}; a silent mismatch would invalidate the
mechanism we argue from.

\paragraph{AST ground truth.} \code{tree-sitter} gives byte-level start/end offsets for
statement-level nodes (\code{if\_statement}, \code{for\_statement}, \code{declaration}) and
expression-level nodes (\code{call\_expression}, \code{binary\_expression}). Separately we
query pointer dereferencing (\code{*}), address referencing (\code{\&}) and manual memory
management (\code{new}/\code{delete}) as memory-unsafe regions.

\paragraph{Adversarial data.} A noised variant of the CodeSearchNet subset injects missing
semicolons, localized indentation changes, camelCase$\leftrightarrow$snake\_case inversions,
and typos in syntax keywords, to probe robustness (O4).

\section{Evaluation Metrics}
\label{sec:metrics}

\noindent
\textbf{Boundary P/R at constituent starts (primary):} precision and recall of
entropy-triggered boundaries against AST node start offsets within tolerance $k$, reported
as a margin over \emph{both} baselines; we expect precision $>$ recall (\S\ref{sec:why}).
\textbf{P/R stratified by node type:} the same metric grouped by constituent type,
separating deterministic from open-ended openers; this tests P4. \textbf{AST boundary IoU
(secondary):} byte-span overlap between patches and AST nodes, expected to understate the
effect per P2 and reported for comparability with prior segmentation work. \textbf{Compute
density (BPP) and entropy heatmaps:} bytes-per-patch profiling and entropy visualization over
source text; we expect high BPP variance in code, more uniform BPP in prose.
\textbf{Degradation shift:} change in the primary metric and in BPP between clean and
adversarial data, indicating robustness relative to subword tokenizers.
\textbf{Memory-specific compute profiling:} BPP in memory-unsafe regions before and after
conditioning on identifier-token entropy, isolating any genuine semantic-risk effect from
the surface-predictability confound.

\section{Feasibility, Risks and Deliverables}

\paragraph{Feasibility and compute.} The study is deliberately cheap. Everything H1 requires
comes from the \emph{patcher}: a 14-layer, 768-dimensional byte-level model of ${\sim}$150M
parameters; the 1B latent transformer is never needed for boundaries. Inference in
\code{bfloat16} fits a single consumer GPU and the analysis is a forward pass, so the
experiment is reproducible on modest hardware; larger corpus sweeps and the adversarial
condition (O4) dispatch as batch jobs on institutional compute if throughput becomes the
limit. \textbf{No model training is required at any stage.}

\paragraph{Risks and mitigations.} \emph{Checkpoint infidelity:} context window and
patch-length distribution are verified against the original configuration \emph{before} any
results are produced; this gates everything downstream. \emph{Newline confound:} the explicit
whitespace baseline (\S\ref{sec:pred}) --- H1 must beat it, not merely H0. \emph{Identifier
entropy swamps syntax:} predicted in advance (\S\ref{sec:why}) and isolated by regression in
O5, so it is reported as a precision loss rather than an unmodelled confound. \emph{Parser
offset drift:} character offsets are round-tripped to UTF-8 byte offsets and validated before
scoring. \emph{Depth incomparable across languages:} alignment is computed per depth and
normalized before cross-language comparison (O3).

\paragraph{Deliverables.} (1)~A boundary-extraction pipeline producing per-byte entropies,
tagged boundaries and BPP statistics for arbitrary source files. (2)~A \code{tree-sitter} AST
mapping engine for Python and C++, including memory-unsafe region queries. (3)~An evaluation
harness implementing P1--P4, both baselines, the permutation test, and depth/node-type
stratification. (4)~Entropy heatmaps and BPP profiles for code and prose. (5)~A final report and
presentation covering all five objectives, including negative results.

\section{Timeline and Division of Work}

\noindent
\emph{Weeks 1--3:} extraction pipeline; \code{tree-sitter} AST mapping for Python/C++;
checkpoint and context-window sanity checks --- a precondition, resolved before any alignment
number is generated. \emph{Weeks 4--6:} H0 and whitespace baselines; P1/P2 alignment on code;
P4 stratification; prose contrast (P3). \emph{Weeks 7--8:} mid-submission write-up;
adversarial dataset; O4. \emph{Weeks 9--11:} memory-unsafe profiling and confound conditioning
(O5); cross-language comparison (O3). \emph{Week 12+:} consolidation, ablations, final
write-up and presentation.

\paragraph{Division of work.} Three parallel streams; analysis, writing and the presentation
are shared. \textbf{Patcher and extraction:} boundary/entropy extraction, checkpoint
fidelity verification, BPP profiling, entropy heatmaps. \textbf{Ground truth and data:}
\code{tree-sitter} integration for both languages, node-type taxonomy, memory-unsafe queries,
adversarial corpus construction. \textbf{Statistics and controls:} H0 and whitespace
baselines, permutation testing, P3 constituency parsing, O5 confound regression.

\begin{thebibliography}{9}
\footnotesize\setlength{\itemsep}{1pt}\setlength{\parskip}{0pt}
\bibitem[Clark et al.(2022)]{clark2022canine} Jonathan H. Clark, Dan Garrette, Iulia Turc,
and John Wieting. 2022. CANINE: Pre-training an efficient tokenization-free encoder for
language representation. \emph{TACL}, 10:73--91.
\bibitem[Feng et al.(2020)]{feng2020codebert} Zhangyin Feng et al. 2020. CodeBERT: A
pre-trained model for programming and natural languages. \emph{Findings of EMNLP 2020}, pages
1536--1547.
\bibitem[Hewitt and Manning(2019)]{hewitt2019structural} John Hewitt and Christopher D.
Manning. 2019. A structural probe for finding syntax in word representations.
\emph{NAACL-HLT}, pages 4129--4138.
\bibitem[Pagnoni et al.(2024)]{pagnoni2024blt} Artidoro Pagnoni et al. 2024. Byte Latent
Transformer: Patches scale better than tokens. \emph{arXiv preprint arXiv:2412.09871}.
\bibitem[Sennrich et al.(2016)]{sennrich2016bpe} Rico Sennrich, Barry Haddow, and Alexandra
Birch. 2016. Neural machine translation of rare words with subword units. \emph{ACL}, pages
1715--1725.
\bibitem[Xue et al.(2022)]{xue2022byt5} Linting Xue et al. 2022. ByT5: Towards a token-free
future with pre-trained byte-to-byte models. \emph{TACL}, 10:291--306.
\end{thebibliography}

\end{document}
